% TOP_PROBLEM: {"id": 30002013, "title": "Broué's Abelian Defect Group Conjecture", "category_id": 17, "category": {"id": 17, "name": "group_theory", "display_name": "Group Theory", "description": "Problems about groups, group actions, representations, and related algebraic structures.", "slug": "group-theory", "order_index": 17, "created_at": "2026-07-31T15:26:25.671Z"}, "status": "open"}
% FIRST_POSTED: 2026-09-27
% !TeX program = lualatex
% ATTEMPT_STATUS: unresolved
\documentclass[11pt]{article}
\usepackage[margin=25mm]{geometry}
\usepackage{fontspec}
\setmainfont{Latin Modern Roman}
\usepackage{amsmath,amssymb,mathtools}
\usepackage{unicode-math}
\setmathfont{Latin Modern Math}
\usepackage{xurl,hyperref}
\hypersetup{colorlinks=true,urlcolor=blue}
\setcounter{secnumdepth}{0}
\setlength{\parindent}{0pt}
\setlength{\parskip}{5pt}
\setlength{\emergencystretch}{3em}
\begin{document}
\section*{\texorpdfstring{Broué's Abelian Defect Group Conjecture}{Broué's Abelian Defect Group Conjecture}}\label{30002013}
\subsection{Definitions and mathematical statement}
Fix a splitting field $k$ of characteristic $p$ for a finite group $G$. A $p$-block algebra is $B=kGe$ for a primitive central idempotent $e$; its defect group $D$ is the associated maximal $p$-subgroup in Brauer's local block theory, defined up to conjugacy. Let $b$ be the Brauer correspondent of $B$ in $kN_G(D)$. If $D$ is abelian, Brou\'e's conjecture asks for a triangulated equivalence
\[
 D^b(B\text{-}\mathrm{mod})\simeq D^b(b\text{-}\mathrm{mod}).
\]
These bounded derived categories consist of bounded complexes of finitely generated modules, with quasi-isomorphisms inverted. The target is an equivalence of derived categories, not merely equality of character counts or a bijection between simple modules. Stronger refinements concerning a splendid equivalence are not imposed by the recorded statement.
\par\smallskip\textit{Formulation and concept references:} \href{http://www.numdam.org/item/AST_1990__181-182__61_0/}{[S1]}.

\subsection{Short English statement}
When a block has an abelian defect group, is its derived representation theory equivalent to that of its corresponding block in the defect group's normalizer?
\subsection{Research attempt: explicit equivalences and necessary derived invariants}
\textbf{Status and the level of the requested conclusion.} This attempt does not establish the equivalence for arbitrary abelian defect groups. The primary source [S1] formulates a categorical assertion, stronger than equality of the numbers of simple modules. A current source check, on 27 September 2026, also found the revised primary manuscript [E1]: its August 2026 version descends specified equivalences for RoCK blocks of double covers of symmetric and alternating groups to a smaller coefficient ring. Neither that restricted result nor its stronger splendid refinement is being substituted for the general target here.

The strategy below is to identify cases where a functor is available, derive an invariant which any proposed functor must preserve, and test whether the invariant could plausibly characterize an equivalence. The last test gives a concrete obstruction to a common shortcut.

\textbf{A complete case: a normal defect group.} If the specified defect group $D$ is normal in $G$, then $N_G(D)=G$. The Brauer correspondent in this identical ambient group is $B$ itself. Thus $b=B$ and the identity functor gives the required triangulated equivalence. This conclusion does not need the additional assumption that $D$ is abelian. In particular, the defect-zero case $D=1$ is already covered.

The observation is elementary, but it also prevents a circular proposed proof. The difficult situation is precisely one in which the normalizer is a proper subgroup and the correspondence of blocks does not provide an equivalence of their module categories by definition. Brauer correspondence identifies particular block idempotents in different algebras; it does not supply the derived functor sought in the conjecture.

\textbf{An explicit matrix-algebra model.} Let $A$ be any finite-dimensional $k$-algebra and $C=M_d(A)$. Put $e=E_{11}$. There is an exact equivalence
\[
 C\text{-}\mathrm{mod}\longrightarrow A\text{-}\mathrm{mod},
 \qquad M\longmapsto eM,
 \tag{337.1}
\]
with inverse $N\mapsto A^d\otimes_A N$, where the left matrix action is on the column $A^d$ and $A$ acts on the right entrywise.

For completeness, the natural map from the reconstructed module to $M$ is
\[
 A^d\otimes_A eM\longrightarrow M,
 \qquad (a_1,\ldots,a_d)^T\otimes m
       \longmapsto\sum_i E_{i1}a_i m.
\]
Its inverse sends $m$ to the column of components $E_{1i}m$, using the identification of each component with $eM$. The identities $\sum_iE_{i1}E_{1i}=1$ and $E_{1i}E_{j1}=\delta_{ij}e$ verify both compositions. The functor $M\mapsto eM$ is exact because applying an idempotent to a short exact sequence preserves exactness. Its inverse is exact because $A^d$ is free as a right $A$-module. Applied degree by degree, these functors preserve cohomology and quasi-isomorphisms, and therefore give an equivalence of bounded derived categories.

For example, if $G=P\times H$, with $P$ an abelian $p$-group and $H$ a $p'$-group, splitting semisimplicity gives
\[
 kH\cong\prod_{\chi}M_{d_\chi}(k),
 \qquad kG\cong\prod_{\chi}M_{d_\chi}(kP).
\]
Each displayed factor is a block, and (337.1) explicitly describes its module and derived categories in terms of $kP$. In this example $P$ is normal, so the required equivalence with the Brauer correspondent is already the identity. The matrix calculation supplies a concrete representation of those categories, rather than merely counting their simple objects.

\textbf{Lemma 337.1 (a necessary numerical condition).} If finite-dimensional algebras $A$ and $C$ have equivalent bounded derived categories of finite-dimensional modules, they have the same number of isomorphism classes of simple modules.

\emph{Proof.} For a finite-length module category its Grothendieck group is free abelian on the classes of the simple modules. Indeed, Jordan--Hölder gives a well-defined composition-multiplicity vector, additive in short exact sequences, and a composition series expresses every module class as the corresponding sum of simple classes. These maps are inverse.

The Grothendieck group of the bounded derived category is the same group. Explicitly, associate to a bounded complex $X$ the Euler class
\[
 [X]\longmapsto\sum_i(-1)^i[H^i(X)].
 \tag{337.2}
\]
A distinguished triangle gives a long exact cohomology sequence; breaking it into short exact sequences shows that the alternating sum in (337.2) is additive. Conversely, the truncation triangles express $[X]$ as $\sum_i(-1)^i[H^i(X)]$ in the derived Grothendieck group. Hence inclusion of modules in degree zero and (337.2) are inverse isomorphisms. A triangulated equivalence induces an isomorphism of these free abelian groups, so their ranks, the numbers of simples, agree. $\square$

This proves a necessary equality, but leaves the morphisms and extension groups almost entirely uncontrolled.

\textbf{A counterexample to reversing the numerical argument.} Take
\[
 A=k,\qquad C=k[\varepsilon]/(\varepsilon^2).
\]
Both algebras have one simple module. They nevertheless have inequivalent bounded derived categories. Let $S=C/(\varepsilon)$. Its projective resolution is
\[
 \cdots\xrightarrow{\varepsilon}C\xrightarrow{\varepsilon}C
 \xrightarrow{\varepsilon}C\longrightarrow S\longrightarrow0.
 \tag{337.3}
\]
At every displayed copy of $C$, the kernel and image of multiplication by $\varepsilon$ are both $k\varepsilon$, so the sequence is exact. Applying $\operatorname{Hom}_C(-,S)$ gives zero differentials, because $\varepsilon$ acts trivially on $S$. Therefore
\[
 \operatorname{Ext}^j_C(S,S)\cong k\qquad(j\geq0).
\]
Equivalently, $\operatorname{Hom}_{D^b(C)}(S,S[j])$ is nonzero for arbitrarily large positive $j$.

Over $k$, every short exact sequence of vector spaces splits. A bounded complex is consequently isomorphic in the derived category to the direct sum of its cohomology spaces placed in their degrees: split cycles, boundaries, and complements in each degree, leaving contractible two-term summands and zero-differential cohomology summands. For any such bounded object $Y$, one has
\[
 \operatorname{Hom}_{D^b(k)}(Y,Y[j])=0
\]
for all sufficiently large $j$, since its finitely many nonzero cohomology degrees no longer overlap. An equivalence taking $S$ to $Y$ would contradict this vanishing. Thus equality of simple-module counts does not even distinguish these elementary categories.

This example is a logical stress test, not a counterexample to the conjecture: these two algebras have not been presented as a block and its Brauer correspondent with the required common defect group. It demonstrates exactly why a proof of a counting conjecture alone cannot imply the requested equivalence.

\textbf{What a constructive approach would have to provide.} One standard route is to build a bounded complex $T$ of finitely generated projective $B$-modules which is a tilting complex: it must have
\[
 \operatorname{Hom}_{K^b(B\text{-}\mathrm{proj})}(T,T[i])=0
 \quad(i\ne0),
\]
its direct summands, shifts, cones, and finite iterated constructions must generate all of $K^b(B\text{-}\mathrm{proj})$, and its endomorphism algebra, with the appropriate opposite convention for left modules, must be $b$. The derived Morita theorem then gives the equivalence. This criterion is a framework, not a complex constructed here for general blocks.

The generation condition cannot be suppressed. For $B=A\times A$, let $T=(A,0)$ in degree zero. It is projective, its shifted self-morphism groups vanish, and its endomorphism algebra is $A^{\mathrm{op}}$. However every object generated from $T$ still has zero second component. It cannot generate $(0,A)$, and $D^b(A\times A)$ cannot be equivalent to $D^b(A)$ when $A\ne0$, by the rank calculation in Lemma 337.1. Thus local self-extension calculations and an attractive endomorphism algebra can coexist with failure of the decisive global condition.

\textbf{Finite diagnostic and boundary of the attempt.} The accompanying exact calculation represents multiplication by $\varepsilon$ in the basis $1,\varepsilon$, verifies that its square is zero and its kernel equals its image, and checks the matrix-unit identities underlying (337.1). The finite initial segment of (337.3) is accompanied by the symbolic periodic argument above; observing several equal ranks alone would not justify an infinite resolution. No computation here claims to produce tilting complexes for unknown blocks.

The attempt proves the target for normal defect groups and gives explicit equivalences for matrix-algebra models, as well as necessary invariants and counterexamples to reversing them. The first missing step for a general abelian defect group is a functor or tilting complex that simultaneously has the correct local endomorphisms and generates the whole derived category. The supplied elementary data do not construct one. The general selected assertion remains unresolved here.

\subsection{Sources}
\begin{itemize}
\item[S1]M. Broue, Asterisque 181-182:61-92, 1990.
\url{http://www.numdam.org/item/AST_1990__181-182__61_0/}.
\textit{Formulation source.}
\item[E1]Yucong Du and Xin Huang, \emph{The refined Brou\'e conjecture for RoCK blocks of double covers of symmetric and alternating groups}, arXiv:2510.02147v6, revised 25 August 2026.
\url{https://arxiv.org/abs/2510.02147v6}.
\textit{Primary source consulted 27 September 2026; the specified RoCK blocks and descent of their known equivalences are the scope used here.}
\end{itemize}
\end{document}
